\documentclass[addpoints,spanish, 12pt,a4paper]{exam}
%\documentclass[answers, spanish, 12pt,a4paper]{exam}
\printanswers
\pointpoints{punto}{puntos}
\hpword{Puntos:}
\vpword{Puntos:}
\htword{Total}
\vtword{Total}
\hsword{Resultado:}
\hqword{Ejercicio:}
\vqword{Ejercicio:}

\usepackage[utf8]{inputenc}
\usepackage[spanish]{babel}
\usepackage{eurosym}
%\usepackage[spanish,es-lcroman, es-tabla, es-noshorthands]{babel}


\usepackage[margin=1in]{geometry}
\usepackage{amsmath,amssymb}
\usepackage{multicol}
\usepackage{yhmath}

\pointsinrightmargin % Para poner las puntuaciones a la derecha. Se puede cambiar. Si se comenta, sale a la izquierda.
\extrawidth{-2.4cm} %Un poquito más de margen por si ponemos textos largos.
\marginpointname{ \emph{\points}}

\usepackage{graphicx}

\graphicspath{{../img/}} 

\newcommand{\class}{2º Bachillerato CCSS}
\newcommand{\examdate}{\today}
\newcommand{\examnum}{Final 3ªEv.}
\newcommand{\tipo}{A}


\newcommand{\timelimit}{90 minutos}

\renewcommand{\solutiontitle}{\noindent\textbf{Solución:}\enspace}


\pagestyle{head}
\firstpageheader{\includegraphics[width=0.2\columnwidth]{header_left}}{\textbf{Departamento de Matemáticas\linebreak \class}\linebreak \examnum}{\includegraphics[width=0.1\columnwidth]{header_right}}
\runningheader{\class}{\examnum}{Página \thepage\ of \numpages}
\runningheadrule


\usepackage{pgf,tikz,pgfplots}
\pgfplotsset{compat=1.15}
\usepackage{mathrsfs}
\usetikzlibrary{arrows}


\begin{document}

\noindent
\begin{tabular*}{\textwidth}{l @{\extracolsep{\fill}} r @{\extracolsep{6pt}} }
\textbf{Nombre:} \makebox[3.5in]{\hrulefill} & \textbf{Fecha:}\makebox[1in]{\hrulefill} \\
 & \\
\textbf{Tiempo: \timelimit} & Tipo: \tipo 
\end{tabular*}
\rule[2ex]{\textwidth}{2pt}
Esta prueba tiene \numquestions\ ejercicios. La puntuación máxima es de \numpoints. 
La nota final de la prueba será la parte proporcional de la puntuación obtenida sobre la puntuación máxima. 

\begin{center}


\addpoints
 %\gradetable[h][questions]
	\pointtable[h][questions]
\end{center}

\noindent
\rule[2ex]{\textwidth}{2pt}

\begin{questions}

% \question[1] Se valorará hasta con un punto el uso de la notación adecuada en los ejercicios propuestos


% 22.9.4
\question Sean A y B dos sucesos asociados a un mismo experimento aleatorio. Suponga que $P(A) = 0. 7$, $P(\overline{B}) = 0.7$ y $P(A\cap B) = 0.2$. 
\begin{parts}
    \part[1] ¿Son A y B independientes? Justifique su respuesta.
    \part[1] Calcule $P(\overline{A} \cap \overline{B})$
\end{parts}



\begin{solution}
Solución: Por costumbre utilizaremos B = Be. 
a) P(A n B) = O, 2 y P(A)P(B) = O, 7(1 — O, 7) = O, 21 P(A n B) # P(A)P(B) A y B no son independientes. 
b) P(A n B) = P(A U B) =1 — P(A U B) = 1 — (P(A) ± P(B) — P(A n B)) =1 — (O, 7 + O, 3 — O, 2) = O, 2  
\end{solution}

\question[1] Un examen de EVAU de Historia consiste en desarrollar 2 de los 4 temas propuestos. Si el temario se compone de 50 temas y un alumno se ha estudiado solo 10. ¿Cuál es la probabilidad de que se sepa al menos uno de los temas propuestos? 
\begin{solution}
    $1-\frac{40}{50}\cdot\frac{39}{49}\cdot\frac{38}{48}\cdot\frac{37}{47}= 0.603169778549718$
\end{solution}

% Problema 21.1.4 (2 puntos) 
\question En un mercado agropecuario el 70\% de las verduras que se comercializan son de proximidad y el resto no. El 30\% de las verduras de proximidad son ecológicas, mientras que de las que no son de proximidad, solo son ecológicas el 10\%. Si un cliente elegido al azar ha realizado una compra de una verdura, calcule las siguientes probabilidades: 
\begin{parts}
    \part[1] Probabilidad de que la verdura comprada no sea ecológica
    \part[1] Probabilidad de que la verdura sea de proximidad o ecológica
\end{parts}
\begin{solution}
$0.76$ y $0.73$  
\end{solution}

% Problema 21.1.5 (2 puntos) 
% \question El número de kilómetros que un corredor entrena a la semana mientras prepara una carrera popular se puede aproximar por una variable aleatoria de distribución normal de media $\mu$ horas y desviación típica $\sigma = 10$ horas
% \begin{parts}
%     \part[1] Se toma una muestra aleatoria simple de 20 atletas, obteniéndose una media muestral de 30 kilómetros. Determine un intervalo de confianza al 95\% para $\mu$
%     \part[2] Suponga que $\mu=28$ kilómetros. Calcule la probabilidad de que al tomar una muestra aleatoria simple de 10 atletas, la media muestral, $\overline{X}$, esté entre 28 y 30 kilómetros
% \end{parts}
% \begin{solution}
%     Solución: 
% N(p; 10) 
% a) n = 20, .TC = 30 y NC = 95% z012 = 1, 96 
% 10 , E= za./2=_ = 1,m = 4, 004 V20 
% IC = (rC — E, Tc + E) = (30 — 4, 383; 30 + 4, 383) = (25, 617; 34, 383) 
% b) p = 28 y rt = 10: 
% 28 — 28 30 — 28 P(28 <7(<30)=P( << a10/ )— 10/ alT  
% P(0 < Z < 0,63) = P(Z < 0,63)— "Z < 0) = 0,7357 — 0.5 = 0,2357 
% \end{solution}

% Problema 20.8.4 (2 puntos) 
\question En un instituto se decide que los alumnos y alumnas solo pueden utilizar un único color (azul o negro) al realizar los exámenes. Dos de cada tres exámenes están escritas en azul. La probabilidad de que un examen escrito en azul sea de una alumna es de 0,7. La probabilidad de que un examen esté escrito en negro y sea de un alumno es 0,2. Se elige un examen al azar. Calcule las siguientes probabilidades:
\begin{parts}
    \part[1] Sea el examen de un alumno
    \part[1] Sabiendo que está escrito en negro, sea de un alumno
\end{parts}
\begin{solution}
    $0.4$ y $0.6$
\end{solution}

% Problema 21.6.5 (2 puntos) 
\question Para que una determinada marca de chocolate estudie entre sus clientes la demanda de sus cajas de bombones, se desea estimar la proporción de cajas grandes en relación al número de cajas de bombones vendidas, P. 
\begin{parts}
    \part[1] Sabiendo que la proporción poblacional de la demanda es P = 0,2, determine el tamaño mínimo necesario de una muestra de ventas de cajas de bombones para garantizar que, con una confianza del 99\%, el margen de error en la estimación no sea mayor que $0.08$
    \begin{solution}
    $E=z_{\alpha/2}\cdot\sqrt{\frac{pq}{n}} \to 0.08=2.5758293035489004\cdot\sqrt{\frac{0.2\cdot0.8}{n}}$
    \\
    $n=2.5758293035489004^2\cdot\frac{0.2\cdot0.8}{0.08^2}= 165.87241502553036 \to n=166$
        
    \end{solution}
    \part[1] Tomada al azar una muestra de 200 cajas de bombones vendidas, se encontró que 25 habían sido cajas grandes. Determine un intervalo de confianza al 95\% para la proporción de cajas grandes en relación a la venta total de cajas de bombones
    \begin{solution}
    $\frac{25}{200}\pm 1.959963984540054\sqrt{\frac{\frac{25}{200}\cdot\frac{175}{200}}{200}}$\\$ (0,07916;0,17084)$
        
    \end{solution}
\end{parts}




\question[2] Se realiza un estudio sobre la calidad de un determinado producto, a través de una encuesta que responde
el 32\% de los encuestados. Si se realizan 1000 encuestas, ¿Qué probabilidad hay de que respondan, al
menos, 300 personas?
\begin{solution}
    $X\sim  \mathcal{B}(1000, 0,32) \to X \approx X' \sim  \mathcal{N}\left(320, \sqrt{\dfrac{32\cdot68}{10}}\right)$ ya que $np=1000\cdot 0.32= 320>5$ y $np=1000\cdot 0.68= 680>5$
    \\
    $P(X\geq 300)\approx P(X'\geq 299,5)=\displaystyle 0.917691634735558$
\end{solution}

% \question Se pretende analizar el consumo anual en alimentación y bebidas en los hogares españoles.
% Dicha variable sigue una distribución normal con una desviación típica de 3.000 euros.
% \begin{parts}
% \part[2] Si deseamos obtener un intervalo de confianza al 96\% para la media de dicha variable
% ¿cuántas familias tenemos que encuestar para que la amplitud del intervalo no sea superior de
% 2.000 euros?
% \part[1] En una muestra de 60 hogares se obtuvo un consumo medio anual en alimentación y
% bebidas de 17.000 euros, halle el intervalo de confianza al 96\% para la media de dicha variable.
% \part[1] Si desde una asociación de consumidores se afirma que «el consumo anual medio en
% alimentación y bebidas en hogares es de 20.000 euros al año». Razone, a la vista del apartado anterior si
% hay motivos para dudar de su afirmación.    
% \end{parts}

\question Se pretende analizar el gasto anual en gasolina de los hogares españoles.
Dicha variable sigue una distribución normal con una desviación típica de 1.000 euros.
\begin{parts}
\part[1] Si deseamos obtener un intervalo de confianza al 95\% para la media de dicha variable
¿a cuántas familias tenemos que encuestar para que la amplitud del intervalo no sea superior de
500 euros?
\begin{solution}
    El error tiene que ser de $500/2=250$. Por tanto, como $E=z_{\alpha/2}\cdot\frac{\sigma}{\sqrt{n}}$ y tenemos los siguientes datos:
    \begin{itemize}
        \item $E=250$
        \item $\sigma=1000$
        \item $P(Z<=z_{\alpha/2}=0.975) \to z_{\alpha/2}=1.959963984540054$
    \end{itemize}
    $\sqrt{n}=\frac{z_{\alpha/2}\cdot\sigma}{E} \to 
    n=\left(\frac{z_{\alpha/2}\cdot\sigma}{E}\right)^2 \to 
    n=\left(\frac{1.959963984540054\cdot1000}{250}\right)^2= 61.463341131106006 \to 
    n = 62$ familias
\end{solution}
\part[1] En una muestra de 50 hogares se obtuvo un gasto medio anual en gasolina de 4.000 euros, halle el intervalo de confianza al 95\% para la media de dicha variable.
\begin{solution}
La media estará  en $\left(\overline{x}+z_{\alpha/2}\cdot\frac{\sigma}{\sqrt{n}}, \overline{x}-z_{\alpha/2}\cdot\frac{\sigma}{\sqrt{n}}\right)
=
\left(4000+1.959963984540054\cdot\frac{1000}{\sqrt{50}}, 4000-1.959963984540054\cdot\frac{1000}{\sqrt{50}}\right)
=
\left( 3722.81923513006, 4277.18076486994 \right)
$    
\end{solution}
\part[1] Si desde una asociación de consumidores se afirma que «el gasto anual medio en
gasolina en hogares es de 3.500 euros al año». Razone, a la vista del apartado anterior si
hay motivos para dudar de su afirmación.
\begin{solution}
    La información es errónea con una probabilidad del 95\% ya que 3500 no está en el intervalo de confianza
\end{solution}
\end{parts}


\end{questions}

    \newgeometry{left=1 cm,bottom=2cm}
% \begin{landscape}
\begin{table}
% \Large
\centering

% \caption{Extracto de tabla de probabilidades de la \textbf{normal estándar $Z(0,1)$}}
\caption{Tabla de probabilidades de la \textbf{normal estándar $Z(0,1)$}}
\label{my-label}

\begin{tabular}{l|llllllllll}
z   & 0       & 0,01    & 0,02    & 0,03    & 0,04    & 0,05    & 0,06    & 0,07    & 0,08    & 0,09    \\
\hline
0   & 0,5     & 0,50399 & 0,50798 & 0,51197 & 0,51595 & 0,51994 & 0,52392 & 0,5279  & 0,53188 & 0,53586 \\
0,1 & 0,53983 & 0,5438  & 0,54776 & 0,55172 & 0,55567 & 0,55962 & 0,56356 & 0,56749 & 0,57142 & 0,57535 \\
0,2 & 0,57926 & 0,58317 & 0,58706 & 0,59095 & 0,59483 & 0,59871 & 0,60257 & 0,60642 & 0,61026 & 0,61409 \\
0,3 & 0,61791 & 0,62172 & 0,62552 & 0,6293  & 0,63307 & 0,63683 & 0,64058 & 0,64431 & 0,64803 & 0,65173 \\
0,4 & 0,65542 & 0,6591  & 0,66276 & 0,6664  & 0,67003 & 0,67364 & 0,67724 & 0,68082 & 0,68439 & 0,68793 \\
0,5 & 0,69146 & 0,69497 & 0,69847 & 0,70194 & 0,7054  & 0,70884 & 0,71226 & 0,71566 & 0,71904 & 0,7224  \\
0,6 & 0,72575 & 0,72907 & 0,73237 & 0,73565 & 0,73891 & 0,74215 & 0,74537 & 0,74857 & 0,75175 & 0,7549  \\
0,7 & 0,75804 & 0,76115 & 0,76424 & 0,7673  & 0,77035 & 0,77337 & 0,77637 & 0,77935 & 0,7823  & 0,78524 \\
0,8 & 0,78814 & 0,79103 & 0,79389 & 0,79673 & 0,79955 & 0,80234 & 0,80511 & 0,80785 & 0,81057 & 0,81327 \\
0,9 & 0,81594 & 0,81859 & 0,82121 & 0,82381 & 0,82639 & 0,82894 & 0,83147 & 0,83398 & 0,83646 & 0,83891 \\
1   & 0,84134 & 0,84375 & 0,84614 & 0,84849 & 0,85083 & 0,85314 & 0,85543 & 0,85769 & 0,85993 & 0,86214 \\
1,1 & 0,86433 & 0,8665  & 0,86864 & 0,87076 & 0,87286 & 0,87493 & 0,87698 & 0,879   & 0,881   & 0,88298 \\
1,2 & 0,88493 & 0,88686 & 0,88877 & 0,89065 & 0,89251 & 0,89435 & 0,89617 & 0,89796 & 0,89973 & 0,90147 \\
1,3 & 0,9032  & 0,9049  & 0,90658 & 0,90824 & 0,90988 & 0,91149 & 0,91309 & 0,91466 & 0,91621 & 0,91774 \\
1,4 & 0,91924 & 0,92073 & 0,9222  & 0,92364 & 0,92507 & 0,92647 & 0,92785 & 0,92922 & 0,93056 & 0,93189 \\
1,5 & 0,93319 & 0,93448 & 0,93574 & 0,93699 & 0,93822 & 0,93943 & 0,94062 & 0,94179 & 0,94295 & 0,94408 \\
1,6 & 0,9452  & 0,9463  & 0,94738 & 0,94845 & 0,9495  & 0,95053 & 0,95154 & 0,95254 & 0,95352 & 0,95449 \\
1,7 & 0,95543 & 0,95637 & 0,95728 & 0,95818 & 0,95907 & 0,95994 & 0,9608  & 0,96164 & 0,96246 & 0,96327 \\
1,8 & 0,96407 & 0,96485 & 0,96562 & 0,96638 & 0,96712 & 0,96784 & 0,96856 & 0,96926 & 0,96995 & 0,97062 \\
1,9 & 0,97128 & 0,97193 & 0,97257 & 0,9732  & 0,97381 & 0,97441 & 0,975   & 0,97558 & 0,97615 & 0,9767  \\
2   & 0,97725 & 0,97778 & 0,97831 & 0,97882 & 0,97932 & 0,97982 & 0,9803  & 0,98077 & 0,98124 & 0,98169 \\
2,1 & 0,98214 & 0,98257 & 0,983   & 0,98341 & 0,98382 & 0,98422 & 0,98461 & 0,985   & 0,98537 & 0,98574 \\
2,2 & 0,9861  & 0,98645 & 0,98679 & 0,98713 & 0,98745 & 0,98778 & 0,98809 & 0,9884  & 0,9887  & 0,98899 \\
2,3 & 0,98928 & 0,98956 & 0,98983 & 0,9901  & 0,99036 & 0,99061 & 0,99086 & 0,99111 & 0,99134 & 0,99158 \\
2,4 & 0,9918  & 0,99202 & 0,99224 & 0,99245 & 0,99266 & 0,99286 & 0,99305 & 0,99324 & 0,99343 & 0,99361 \\
2,5 & 0,99379 & 0,99396 & 0,99413 & 0,9943  & 0,99446 & 0,99461 & 0,99477 & 0,99492 & 0,99506 & 0,9952  \\
2,6 & 0,99534 & 0,99547 & 0,9956  & 0,99573 & 0,99585 & 0,99598 & 0,99609 & 0,99621 & 0,99632 & 0,99643 \\
2,7 & 0,99653 & 0,99664 & 0,99674 & 0,99683 & 0,99693 & 0,99702 & 0,99711 & 0,9972  & 0,99728 & 0,99736 \\
2,8 & 0,99744 & 0,99752 & 0,9976  & 0,99767 & 0,99774 & 0,99781 & 0,99788 & 0,99795 & 0,99801 & 0,99807 \\
2,9 & 0,99813 & 0,99819 & 0,99825 & 0,99831 & 0,99836 & 0,99841 & 0,99846 & 0,99851 & 0,99856 & 0,99861 \\
3   & 0,99865 & 0,99869 & 0,99874 & 0,99878 & 0,99882 & 0,99886 & 0,99889 & 0,99893 & 0,99896 & 0,999   \\
3,1 & 0,99903 & 0,99906 & 0,9991  & 0,99913 & 0,99916 & 0,99918 & 0,99921 & 0,99924 & 0,99926 & 0,99929 \\
3,2 & 0,99931 & 0,99934 & 0,99936 & 0,99938 & 0,9994  & 0,99942 & 0,99944 & 0,99946 & 0,99948 & 0,9995  \\
3,3 & 0,99952 & 0,99953 & 0,99955 & 0,99957 & 0,99958 & 0,9996  & 0,99961 & 0,99962 & 0,99964 & 0,99965 \\
3,4 & 0,99966 & 0,99968 & 0,99969 & 0,9997  & 0,99971 & 0,99972 & 0,99973 & 0,99974 & 0,99975 & 0,99976 \\
3,5 & 0,99977 & 0,99978 & 0,99978 & 0,99979 & 0,9998  & 0,99981 & 0,99981 & 0,99982 & 0,99983 & 0,99983 \\
3,6 & 0,99984 & 0,99985 & 0,99985 & 0,99986 & 0,99986 & 0,99987 & 0,99987 & 0,99988 & 0,99988 & 0,99989 \\
3,7 & 0,99989 & 0,9999  & 0,9999  & 0,9999  & 0,99991 & 0,99991 & 0,99992 & 0,99992 & 0,99992 & 0,99992 \\
3,8 & 0,99993 & 0,99993 & 0,99993 & 0,99994 & 0,99994 & 0,99994 & 0,99994 & 0,99995 & 0,99995 & 0,99995 \\
3,9 & 0,99995 & 0,99995 & 0,99996 & 0,99996 & 0,99996 & 0,99996 & 0,99996 & 0,99996 & 0,99997 & 0,99997 \\
4   & 0,99997 & 0,99997 & 0,99997 & 0,99997 & 0,99997 & 0,99997 & 0,99998 & 0,99998 & 0,99998 & 0,99998
\end{tabular}
\end{table}
% \end{landscape}
\restoregeometry

\end{document}
\grid
